\documentclass[11pt,a4paper]{article}
\usepackage[margin=2.3cm]{geometry}
\usepackage{enumitem}
\usepackage{titlesec}
\usepackage{hyperref}
\usepackage{xcolor}
\usepackage{booktabs}
\usepackage{parskip}
\usepackage{amsmath}
\usepackage{float}
\usepackage{longtable}
\usepackage{listings}
\usepackage{tcolorbox}
\tcbuselibrary{listings,breakable,skins}

\definecolor{accent}{RGB}{30,70,120}
\definecolor{green}{RGB}{40,120,40}
\definecolor{red}{RGB}{180,40,40}
\definecolor{codebg}{RGB}{245,245,245}
\definecolor{codeframe}{RGB}{200,200,200}
\definecolor{codekw}{RGB}{0,100,170}
\definecolor{codecomment}{RGB}{100,130,100}
\definecolor{codestr}{RGB}{170,55,55}
\definecolor{strongbg}{RGB}{230,245,230}
\definecolor{strongframe}{RGB}{60,140,60}
\definecolor{warnbg}{RGB}{255,250,230}
\definecolor{warnframe}{RGB}{180,150,40}
\definecolor{failbg}{RGB}{255,235,235}
\definecolor{failframe}{RGB}{180,60,60}
\definecolor{infobg}{RGB}{235,242,250}
\definecolor{infoframe}{RGB}{70,110,170}
\definecolor{kimbg}{RGB}{243,236,250}
\definecolor{kimframe}{RGB}{120,70,170}

\newtcolorbox{resultbox}[2][]{
    enhanced, colback=#2bg, colframe=#2frame,
    boxrule=0.8pt, arc=3pt, left=6pt, right=6pt, top=4pt, bottom=4pt,
    fonttitle=\bfseries, coltitle=white, title={#1}, breakable,
}
\newtcolorbox{kimbox}[1]{
    enhanced, colback=kimbg, colframe=kimframe,
    boxrule=0.9pt, arc=3pt, left=6pt, right=6pt, top=4pt, bottom=4pt,
    fonttitle=\bfseries, coltitle=white, title={#1}, breakable,
}

\lstdefinestyle{shell}{
    basicstyle=\ttfamily\tiny, keywordstyle=\color{codekw},
    numbers=left, numberstyle=\tiny\color{gray}, numbersep=7pt,
    frame=leftline, framesep=5pt, rulecolor=\color{codeframe},
    commentstyle=\color{codecomment}\itshape, stringstyle=\color{codestr},
    showstringspaces=false, breaklines=true, breakatwhitespace=false,
    columns=fullflexible, tabsize=2, numbers=none,
    literate={\$}{{\textcolor{codekw}{\$}}}1,
}
\lstdefinestyle{py}{
    language=Python, basicstyle=\ttfamily\tiny,
    numbers=left, numberstyle=\tiny\color{gray}, numbersep=7pt,
    frame=leftline, framesep=5pt, rulecolor=\color{codeframe},
    keywordstyle=\color{codekw}\bfseries,
    commentstyle=\color{codecomment}\itshape, stringstyle=\color{codestr},
    showstringspaces=false, breaklines=true, columns=fullflexible, tabsize=4,
    numbers=none,
}
\newtcolorbox{codebox}[1][]{
    colback=codebg, colframe=codeframe, boxrule=0.5pt, arc=2pt,
    left=2pt, right=2pt, top=2pt, bottom=2pt, breakable,
    title={\small\ttfamily #1}, fonttitle=\bfseries\color{accent},
}

\titleformat{\section}{\large\bfseries\color{accent}}{\thesection.\ }{0.4em}{}[\titlerule]
\titleformat{\subsection}{\normalsize\bfseries\color{accent}}{\thesubsection\ }{0.4em}{}
\titleformat{\part}[display]{\Huge\bfseries\color{accent}}{\partname\ \thepart}{0.3em}{\titlerule\vspace{0.4em}}

\hypersetup{colorlinks=true, linkcolor=accent, urlcolor=accent}

\begin{document}

\begin{center}
\vspace*{1.2cm}
{\Huge\bfseries\color{accent} Going Live}\\[1em]
{\LARGE Deploying the Copper--Gas--Bitcoin--Sol\\[0.3em]
Recycling Portfolio on Oracle Cloud}\\[1.8em]
{\large A complete runbook}\\[0.8em]
{\small Robin Langer\quad·\quad August 2026}\\[2.5em]
\begin{minipage}{0.84\textwidth}
\itshape\centering
This document is written to be executed. Claude reads it and does the
engineering. Kim does the four things that legally and physically
cannot be delegated: proving who he is, funding the accounts, and
authorising the keys.
\end{minipage}
\vspace*{1cm}
\end{center}
\thispagestyle{empty}
\newpage

\tableofcontents
\newpage

% ══════════════════════════════════════════════════════════════════
\part{Orientation}

\section{Who Does What}

There are three actors, and the boundaries matter --- both for security
and because some steps are legally personal to the account holder.

\vspace{0.5em}
\begin{tabular}{@{}p{2.4cm}p{11.5cm}@{}}
    \toprule
    \textbf{Actor} & \textbf{Responsibilities} \\
    \midrule
    \textbf{Kim} &
    Identity verification (KYC) at each venue; funding accounts from his
    own bank; generating API keys and choosing their permissions;
    approving anything that moves money out. \emph{Cannot be delegated.}
    \\[0.4em]
    \textbf{Claude} &
    Everything technical: provisioning the server, SSH, installing the
    scripts, the virtualenv, cron, the signal engine, the execution
    layer, the weekly reports, monitoring, and debugging.
    \\[0.4em]
    \textbf{Robin} &
    Owns the strategy and the research. Reviews the go-live decision at
    each phase gate. Holds nothing of Kim's.
    \\
    \bottomrule
\end{tabular}

\vspace{1em}
\begin{resultbox}[The security model, stated once and applied throughout]{fail}
\textbf{Kim never hands over a bank login, a password, or a
seed phrase.} Not to Robin, not to Claude, not to anyone. Anybody who
asks for those is either mistaken or stealing.

\vspace{0.5em}
What Kim \emph{does} hand over is a set of \textbf{API keys scoped to
trading only, with withdrawal permission switched off and locked to
the server's IP address}. That is the entire difference between
``a system that can rebalance my portfolio'' and ``a system that can
empty my account.'' An API key without withdrawal rights cannot move
money off the venue --- the worst a total compromise achieves is bad
trades, which is bounded, recoverable, and insurable by position
limits.

\vspace{0.5em}
Kim's banking details are typed by Kim, into the broker's own deposit
screen, and go nowhere else. See Section~\ref{sec:banking}.
\end{resultbox}

\section{What Already Exists}

Not starting from zero. As of August 2026:

\begin{itemize}[leftmargin=2em]
    \item \textbf{An OCI instance is already running} ---
          \texttt{ubuntu@207.211.145.30}, Ubuntu 24.04.4 LTS on
          \texttt{aarch64} (ARM), Python 3.12.3, 23~GB RAM, 69~GB free
          disk, \texttt{cron} active with ten existing jobs. SSH alias
          \texttt{oci}. It hosts four AI agents in Docker plus the
          Gateway app.
    \item \textbf{The research is done} --- see
          \texttt{algorithm.pdf} (the complete, self-contained code that
          generates the portfolio comparison),
          \texttt{kim.pdf} (the full cross-asset study) and
          \texttt{concepts.pdf} (the primer for Kim).
    \item \textbf{What does not exist yet} --- the execution layer (code
          that places orders), any funded account, and a measurement of
          how much of the copper and gas alpha survives in a tradeable
          instrument. All three are addressed below.
\end{itemize}

\noindent
Note: \texttt{\textasciitilde/git/oci} on Robin's Mac is a plain
directory, \emph{not} a git repository, and there is no GitHub repo
named \texttt{oci} under either \texttt{RaggedR} or
\texttt{nickmeinhold}. If you want this versioned, it needs
\texttt{git init} and a new remote.

\section{Where It Runs, and the SSH Question}

The reports run on \textbf{Robin's existing instance}
(\texttt{207.211.145.30}). It is already provisioned, already has cron,
and costs nothing. Kim needs no cloud account at all.

Kim runs nothing on the server himself. He works from \textbf{his own
machine}, where he does exactly two things: fetch the weekly report, and
--- optionally --- log in to look around.

\subsection{No, you do not hand over your SSH credentials}

The question is the right one to ask, and the answer is unambiguous:

\begin{resultbox}[Never share a private key. With anyone. For any reason.]{fail}
An SSH private key is not a password that can be scoped or revoked
per-person --- it \emph{is} your identity on every machine that trusts
it. Sharing \texttt{\textasciitilde/.ssh/oci\_new} would give Kim
Robin's access to Robin's server, indistinguishably from Robin, with no
way to tell the two apart in a log and no way to revoke one without
revoking both.

\vspace{0.5em}
That is not how SSH is meant to be used, and it is not necessary.
\textbf{Public-key authentication is designed for exactly this
situation.}
\end{resultbox}

\subsection{How it actually works}

Keys come in pairs. The \textbf{private} half never leaves the machine
that will connect. The \textbf{public} half is meant to be handed out
--- that is its entire purpose --- and it is useless for anything except
recognising its private twin.

So: \textbf{Kim generates his own pair on his own laptop and sends Robin
only the public half.} Robin adds that line to the server. Now Kim logs
in as Kim, with a key Robin has never seen, and Robin's own key is
untouched.

\vspace{0.5em}
\begin{tabular}{@{}p{3.4cm}p{10.6cm}@{}}
    \toprule
    \textbf{Property} & \textbf{Consequence} \\
    \midrule
    Robin's key unchanged & Nothing to rotate, nothing exposed. \\
    Kim has his own identity & The logs distinguish them. \\
    Revocation is one line & Delete Kim's line from
    \texttt{authorized\_keys} and his access ends instantly, without
    affecting Robin. \\
    Separate Unix account & Kim's files, keys and reports live under
    \texttt{/home/kim}, not mixed into Robin's home directory. \\
    \bottomrule
\end{tabular}

\subsection{Kim's side: generate a key}

\begin{kimbox}{Step 1: make an SSH key --- 5 minutes}
\begin{enumerate}[leftmargin=1.6em, itemsep=3pt]
    \item Open Terminal (macOS/Linux), or PowerShell on Windows 10+.
    \item Run, exactly as written:

    \vspace{0.3em}
    \texttt{ssh-keygen -t ed25519 -C "kim-portfolio"}

    \vspace{0.3em}
    Press Enter to accept the default location. \textbf{Set a
    passphrase} when asked --- if your laptop is ever stolen, that
    passphrase is the only thing between the thief and the server.

    \item Print the \textbf{public} half:

    \vspace{0.3em}
    \texttt{cat \textasciitilde/.ssh/id\_ed25519.pub}

    \vspace{0.3em}
    \item Send Robin \textbf{that whole line}. It starts with
          \texttt{ssh-ed25519} and ends with \texttt{kim-portfolio}.
          Email is fine for this one --- it is a public key, and it is
          safe by design.

    \item \textbf{Never send the other file}
          (\texttt{\textasciitilde/.ssh/id\_ed25519}, no
          \texttt{.pub}). That one is yours forever. If anyone ever
          asks you for it, the answer is no.
\end{enumerate}
\end{kimbox}

\subsection{Robin's side: add Kim, in one command}

\begin{codebox}[Run once, from Robin's Mac]
\begin{lstlisting}[style=shell]
KIMKEY='ssh-ed25519 AAAA...  kim-portfolio'     # the line Kim sent

ssh oci "sudo bash -s" <<EOF
set -euo pipefail
# A normal login user. Deliberately NOT added to the sudo group.
id kim >/dev/null 2>&1 || sudo adduser --disabled-password --gecos "" kim
install -d -m 700 -o kim -g kim /home/kim/.ssh
touch /home/kim/.ssh/authorized_keys
grep -qF "kim-portfolio" /home/kim/.ssh/authorized_keys \
  || echo "$KIMKEY" >> /home/kim/.ssh/authorized_keys
chmod 600 /home/kim/.ssh/authorized_keys
chown -R kim:kim /home/kim/.ssh
echo "kim added; sudo group membership:"; groups kim
EOF
\end{lstlisting}
\end{codebox}

\begin{codebox}[Kim then adds this to his own ~/.ssh/config]
\begin{lstlisting}[style=shell]
Host portfolio
  HostName 207.211.145.30
  User kim
  IdentityFile ~/.ssh/id_ed25519
\end{lstlisting}
\end{codebox}

\noindent
and \texttt{ssh portfolio} works from Kim's laptop. To revoke, ever:

\begin{codebox}[Revocation]
\begin{lstlisting}[style=shell]
ssh oci "sudo sed -i '/kim-portfolio/d' /home/kim/.ssh/authorized_keys"
# or remove the account entirely:
ssh oci "sudo deluser --remove-home kim"
\end{lstlisting}
\end{codebox}

\subsection{The trust boundary, stated honestly}

Two things follow from sharing a machine, and both should be said out
loud rather than discovered later.

\begin{enumerate}[leftmargin=2em]
    \item \textbf{Robin has root. Robin can read Kim's API keys.}
          No file permission changes this; that is what root means.
          The arrangement rests on trust between brothers, which is
          fine --- but it is trust, not a technical guarantee, and Kim
          should know which one he is relying on.
    \item \textbf{Kim gets no \texttt{sudo}.} Not as a slight: the box
          also runs four AI agents and the Gateway app, and an
          accidental \texttt{sudo} command from someone unfamiliar with
          it could take those down. Kim needs nothing more than read
          access to his own report directory.
\end{enumerate}

\vspace{0.5em}
\begin{resultbox}[What makes this acceptable]{strong}
The reason a shared box is tolerable here is the design decision in
Section~\ref{sec:keys}: \textbf{the API keys have withdrawal disabled.}
The worst outcome from any compromise --- Robin's, Kim's, or an
agent's --- is unwanted trades, not stolen funds.

\vspace{0.4em}
If that ever stops being true, or if the amount grows past the point
where ``I trust my brother's server'' is a sentence Kim is comfortable
saying, the answer is a separate Always Free tenancy in Kim's own name.
That is an hour of clicking, the scripts are host-portable, and
\texttt{install\_oci.sh} takes an \texttt{OCI\_HOST} variable
specifically so the move is a one-line change.
\end{resultbox}

\subsection{Kim's laptop: fetching reports}

\begin{kimbox}{Step 2: get the weekly report --- 2 minutes, once}
\begin{enumerate}[leftmargin=1.6em, itemsep=3pt]
    \item The report is written every Monday morning whether or not
          anyone looks. Nothing breaks if you ignore it for a month.
    \item The numbers, without downloading anything:

          \vspace{0.3em}
          \texttt{ssh portfolio 'cat \textasciitilde/portfolio/reports/latest.json'}

          \vspace{0.3em}
    \item The PDF, copied to your Desktop:

          \vspace{0.3em}
          \texttt{scp portfolio:\textasciitilde/portfolio/reports/latest.pdf \textasciitilde/Desktop/}

          \vspace{0.3em}
    \item If you would rather not touch a terminal at all, say so ---
          Robin can have the server email or message you the PDF on
          Monday instead. There is no requirement for Kim to use SSH.
\end{enumerate}
\end{kimbox}


\section{Step 2 --- Crypto Exchange for Bitcoin and Solana}
\label{sec:crypto}

\textbf{Why an exchange and not just a wallet:} this is the single most
common point of confusion, and Section~\ref{sec:wallets} explains it
properly. The short version: a \emph{wallet} holds coins; an
\emph{exchange account} can buy and sell them. A strategy that trades
needs an exchange account with an API. We will use a wallet too, later,
as a vault.

\textbf{Recommended venue: Kraken.} Long operating history, serves
Australian customers with AUD deposits, lists both Bitcoin and Solana,
and --- decisively for us --- has a straightforward API-key system with
\emph{granular permissions} and no requirement to re-authenticate by
hand every day. Verify at signup that it still serves your state; if
not, the alternatives are Independent Reserve, BTC Markets, or Coinbase
Advanced, all of which have comparable APIs.

\begin{kimbox}{Step 2: open and verify a Kraken account --- 30 minutes plus 1--3 days verification}
\begin{enumerate}[leftmargin=1.6em, itemsep=3pt]
    \item Go to \url{https://www.kraken.com} and create an account with
          a \textbf{unique password from your password manager}.
    \item Turn on two-factor authentication immediately, using your
          authenticator app. \textbf{Not} SMS --- SIM-swap attacks on
          crypto accounts are routine.
    \item Complete identity verification: photo ID, a selfie, and your
          address. Choose the tier that permits AUD deposits and
          trading. Approval usually takes hours; occasionally days.
    \item \textbf{Do not deposit anything yet.} Funding happens at
          Phase 3, after the system has been proven on paper.
    \item Now create the API key. \textbf{Settings $\rightarrow$ API
          $\rightarrow$ Add key}, and set permissions
          \emph{exactly} as follows:

          \vspace{0.4em}
          \begin{tabular}{@{}ll@{}}
              \toprule
              \textbf{Permission} & \textbf{Setting} \\
              \midrule
              Query funds / balances       & \textcolor{green}{\textbf{ON}} \\
              Query open \& closed orders  & \textcolor{green}{\textbf{ON}} \\
              Create \& modify orders      & \textcolor{green}{\textbf{ON}} \\
              Cancel / close orders        & \textcolor{green}{\textbf{ON}} \\
              \textbf{Withdraw funds}      & \textcolor{red}{\textbf{OFF --- critical}} \\
              Deposit funds                & \textcolor{red}{\textbf{OFF}} \\
              Access WebSockets            & \textcolor{green}{\textbf{ON}} \\
              Export data                  & \textcolor{green}{\textbf{ON}} \\
              \bottomrule
          \end{tabular}

          \vspace{0.4em}
    \item If Kraken offers to restrict the key to specific IP
          addresses, \textbf{say yes} and enter the server address
          Robin gives you. This means a stolen key is useless from
          anywhere else on earth.
    \item Kraken shows the \textbf{API Secret exactly once.} Copy both
          the Key and the Secret straight into your password manager
          before closing the page.
    \item \textbf{Send Robin} the API Key and Secret --- and send them
          \textbf{through a channel that is not email}: Signal, or a
          password-manager shared item. Email is stored forever on
          servers you do not control.
\end{enumerate}
\end{kimbox}

\begin{resultbox}[Why ``Withdraw funds: OFF'' is the most important switch in this document]{strong}
With withdrawal off, the absolute worst case --- server compromised,
keys stolen, attacker in full control --- is that somebody makes
stupid trades inside Kim's account. Unpleasant, bounded, and partly
recoverable.

With withdrawal on, the worst case is that the money is gone in
ninety seconds and is never coming back. Crypto transactions are
irreversible by design; there is no chargeback and no fraud
department.

The convenience gained by enabling withdrawal is zero, because
nothing in this system ever needs to move money off the exchange.
\end{resultbox}

\section{Step 3 --- Broker for Copper and Natural Gas}
\label{sec:broker}

Copper and natural gas are not coins; they are commodities, traded as
futures contracts on the Chicago Mercantile Exchange. This leg needs a
real broker, and there is a genuine problem to solve first.

\begin{resultbox}[The problem: the contracts are far too big]{fail}
Our backtest used front-month futures prices (\texttt{HG=F},
\texttt{NG=F}). Buying one of those contracts commits you to a great
deal more money than the whole portfolio:

\vspace{0.5em}
\centering
\begin{tabular}{@{}llr@{}}
    \toprule
    \textbf{Contract} & \textbf{Size} & \textbf{Approx.\ notional} \\
    \midrule
    Copper (HG)             & 25{,}000 lb    & $\sim$\$110{,}000 \\
    Micro copper (MHG)      & 2{,}500 lb     & $\sim$\$11{,}000 \\
    Henry Hub natgas (NG)   & 10{,}000 MMBtu & $\sim$\$29{,}000 \\
    E-mini natgas (QG)      & 2{,}500 MMBtu  & $\sim$\$7{,}250 \\
    \bottomrule
\end{tabular}

\vspace{0.5em}
\raggedright
On a \$10{,}000 portfolio, a single copper contract is eleven times
the account. Even the micro contract is more than the entire intended
copper allocation. Futures are leveraged instruments --- you post a
margin deposit, not the full value --- which means a position that
looks affordable can lose more than you deposited. \textbf{The
backtest assumed no leverage. Trading these contracts directly would
not be implementing the strategy; it would be implementing a much
riskier strategy that happens to share its signals.}
\end{resultbox}

\subsection{The alternative, and its honest cost}

The way to get unleveraged, correctly-sized exposure is an
exchange-traded fund that tracks the commodity: \textbf{CPER} for
copper, \textbf{UNG} for natural gas. These trade like shares, can be
bought in any quantity (fractionally, at some brokers), and cannot lose
more than you put in.

The cost of that convenience is real and must be stated:

\begin{enumerate}[leftmargin=2em]
    \item \textbf{Roll decay.} These funds hold futures and must
          repeatedly sell expiring contracts to buy later ones. When
          later contracts are more expensive --- the usual state of
          affairs in natural gas --- this bleeds value continuously.
          \textbf{UNG has historically underperformed the natural gas
          price it tracks by a wide margin.} Our backtest measured
          \texttt{NG=F} and will therefore \emph{overstate} what UNG
          delivers. The strategy's saving grace is that it holds gas
          only about half the time, but the drag is not eliminated,
          merely reduced.
    \item \textbf{Tax complexity.} Both are structured as commodity
          pools and issue US K-1 style tax documents rather than
          ordinary dividend statements. For an Australian taxpayer this
          is a genuine accounting nuisance. Ask an accountant before
          Phase 3.
    \item \textbf{Not all brokers permit them.} Verify at signup.
\end{enumerate}

\begin{resultbox}[Recommendation, stated plainly]{warn}
\textbf{Trade the crypto leg for real; keep the copper and gas leg on
paper until the instrument question is settled.}

The Bitcoin and Solana legs are clean: spot, unleveraged, one venue,
a good API. They can go live as soon as Kim is verified and funded.

The copper and gas legs have an unresolved instrument-selection
problem --- futures are too large, and the ETFs carry decay and tax
costs the backtest never measured. Running them on paper for a quarter
costs nothing and tells us how much of the backtested alpha actually
survives in the tradeable instrument. \textbf{That measurement should
be made before Kim's money is committed, not after.}
\end{resultbox}

\subsection{Which broker}

\textbf{Position: Interactive Brokers.} It is the standard answer ---
access to CME futures, US ETFs, an Australian entity, and a
well-documented API that \texttt{ib\_async} speaks fluently.

\textbf{Objection.} IBKR's API is built around a desktop Java
application (Trader Workstation or IB Gateway) that must be running for
the API to answer, and its session requires two-factor
re-authentication roughly daily. On a headless ARM server this means a
Java gateway under a virtual framebuffer, plus Kim tapping ``approve''
on his phone most days. For a system whose entire selling point is that
it runs unattended, a daily human tap is a design flaw, not a detail.

\textbf{Resolution.} Choose the broker by what the automation needs,
not by reputation. For the ETF route --- which is what we recommend
--- a broker with plain API-key authentication and no daily re-auth is
strictly better. \textbf{Alpaca} is the natural fit: pure REST, keys
that do not expire daily, fractional shares, a built-in paper-trading
environment that is API-identical to live, and support for
US-listed ETFs. Keep IBKR in reserve for the day we decide we
genuinely want futures exposure and accept the operational burden.

\textbf{Practical difference:} Alpaca can be automated tonight, in
paper mode, with no account at all. IBKR cannot be automated unattended
without solving the 2FA problem first.

\begin{kimbox}{Step 3: open a broker account --- 30 minutes plus verification}
\begin{enumerate}[leftmargin=1.6em, itemsep=3pt]
    \item \textbf{Wait.} This step is deliberately last and is not
          needed for Phases 1--2. Do not open it until Robin confirms
          the instrument decision.
    \item When told to proceed, go to \url{https://alpaca.markets},
          create an account, and confirm during signup that Australian
          residents are accepted and that \texttt{CPER} and
          \texttt{UNG} are tradeable. If either answer is no, tell
          Robin and stop --- do not improvise a substitute instrument.
    \item Complete identity verification as in Step 2. Enable
          two-factor authentication.
    \item Generate API keys for the \textbf{paper} environment first
          and send those to Robin. Paper keys cannot touch money, so
          they are safe to hand over early and they let Claude prove
          the whole pipeline works.
    \item Generate \textbf{live} keys only when Robin confirms
          Phase~3. Restrict to trading; no funds-transfer permission;
          IP-restricted if offered.
    \item \textbf{Send Robin} the key and secret over Signal, never
          email. Label clearly which are paper and which are live.
\end{enumerate}
\end{kimbox}

\section{Step 4 --- Funding, and Your Banking Details}
\label{sec:banking}

\begin{kimbox}{Step 4: funding --- only at Phase 3, and only by you}
\begin{enumerate}[leftmargin=1.6em, itemsep=3pt]
    \item Decide the amount, and apply one test: \textbf{if this money
          vanished entirely, would anything in your life break?} If
          yes, the number is too big. Read the drawdown table in
          \texttt{concepts.pdf} before answering --- an $80\%$
          temporary loss is inside the historical range of what this
          strategy does.
    \item Log in to the venue \textbf{yourself}. Find Deposit. Choose
          AUD bank transfer (PayID/Osko is usually instant; a normal
          transfer takes a day).
    \item \textbf{Type your own bank details into the venue's deposit
          screen.} Copy the reference code exactly --- deposits
          without it get held up.
    \item Transfer from your bank as you would to any payee. The first
          transfer may be held for review; this is routine.
    \item \textbf{Tell Robin the amount and the date. Send no account
          numbers, no BSB, no bank password, no card number.} Nothing
          in this system needs them, and nobody who does need them
          would ask you through a person.
    \item Start with the smallest amount that makes the exercise real
          --- a few hundred dollars, not the full allocation. Phase 4
          is where the size increases, after Phase 3 has run clean.
\end{enumerate}
\end{kimbox}

\begin{resultbox}[A note on whose money this is]{warn}
The accounts are in \textbf{Kim's own name}, funded from Kim's own
bank, and Kim can close them or revoke the keys at any moment without
asking anyone. That is deliberate, and it should stay that way.

It keeps custody with Kim, and it keeps the arrangement clearly on the
right side of a real regulatory line: in Australia, managing other
people's money as a business generally requires an Australian
Financial Services licence. A brother sharing research and
infrastructure is not that. Pooled funds, a fee, or discretion over
someone else's account would be. If this ever grows past two brothers
and a shared spreadsheet, get advice first.
\end{resultbox}

% ══════════════════════════════════════════════════════════════════
\part{Wallets, Accounts, and Keys}

\section{Wallet Versus Exchange Account}
\label{sec:wallets}

Four different things get called ``a wallet,'' and conflating them is
how people lose money.

\vspace{0.5em}
\begin{tabular}{@{}p{3.0cm}p{5.2cm}p{5.4cm}@{}}
    \toprule
    \textbf{Thing} & \textbf{What it is} & \textbf{Can it trade?} \\
    \midrule
    Exchange account &
    An account at a company (Kraken) that holds coins on your behalf and
    matches your buy and sell orders. &
    \textbf{Yes} --- and it is the only one of these four that can.
    This is what the strategy uses. \\[0.4em]
    Software wallet &
    An app holding your private keys (Phantom for Solana, Sparrow for
    Bitcoin). You control the coins directly. &
    No. It can send and receive, not trade. \\[0.4em]
    Hardware wallet &
    A dedicated device (Ledger, Trezor) keeping keys offline. The
    gold standard for storage. &
    No. \\[0.4em]
    Paper backup &
    Twelve or twenty-four words written down. \emph{Is} the wallet, in
    every sense that matters. &
    No. \\
    \bottomrule
\end{tabular}

\vspace{0.5em}

\begin{resultbox}[The one rule about seed phrases]{fail}
A seed phrase \textbf{is} the money. Anyone holding those words owns
the coins, permanently and irreversibly.

\vspace{0.4em}
Therefore: write it on paper, never photograph it, never type it into
anything except the wallet app that generated it, never put it in
email or cloud storage, and \textbf{never send it to Robin, to Claude,
or to any person or program that asks.} There is no legitimate reason
for anyone to see it. Every message that has ever asked for one was a
theft in progress.
\end{resultbox}

\subsection{So do we need a wallet at all?}

For trading, no. For safety, yes --- eventually.

Leaving coins on an exchange means trusting the exchange to stay
solvent and unhacked, and the history of the industry is unkind on both
counts. The rule of thumb: the exchange holds the working capital the
strategy is actively trading; a hardware wallet holds anything you have
decided to keep.

\begin{kimbox}{Optional Step 5: a hardware wallet --- do this at Phase 4, not now}
\begin{enumerate}[leftmargin=1.6em, itemsep=3pt]
    \item Buy a Ledger or Trezor \textbf{direct from the
          manufacturer's own website}. Never from eBay, Amazon
          marketplace, or a third-party seller --- tampered devices with
          pre-generated seed phrases are a known and effective attack.
    \item Initialise it yourself. It will show you a seed phrase.
          Write it on paper. Store it somewhere a burglar and a house
          fire would not both reach.
    \item Bitcoin: use the device's built-in Bitcoin app.
          Solana: use the device together with Phantom
          (\url{https://phantom.app}), which supports hardware signing.
    \item Send a \textbf{small test amount} first --- \$20 --- and
          confirm it arrives before moving anything meaningful.
          Crypto addresses have no ``are you sure?'' and no recall.
    \item This wallet is a vault, not part of the trading system. The
          scripts never see it and never need to.
\end{enumerate}
\end{kimbox}

\section{Where Keys Live on the Server}
\label{sec:keys}

Keys go in one file, readable only by the owning user, and are read from
the environment. Never in the scripts, never in git.

\begin{codebox}[Claude: create /home/ubuntu/portfolio/.env]
\begin{lstlisting}[style=shell]
# On the server. Note the permissions -- this matters.
install -m 700 -d /home/ubuntu/portfolio
cat > /home/ubuntu/portfolio/.env <<'EOF'
# Kraken -- trade permissions only, withdrawal DISABLED
KRAKEN_API_KEY=...
KRAKEN_API_SECRET=...

# Alpaca -- paper endpoint until Phase 3
ALPACA_API_KEY=...
ALPACA_API_SECRET=...
ALPACA_BASE_URL=https://paper-api.alpaca.markets

# Risk limits, enforced in code (see section on the execution layer)
MAX_ORDER_USD=500
SOL_CAP=0.25
LIVE_TRADING=0
EOF
chmod 600 /home/ubuntu/portfolio/.env
\end{lstlisting}
\end{codebox}

\begin{resultbox}[Rules for handling Kim's keys]{fail}
\begin{enumerate}[leftmargin=1.5em, topsep=2pt]
    \item \texttt{chmod 600}. Never world-readable.
    \item Never committed to git. Add \texttt{.env} to
          \texttt{.gitignore} \emph{before} the first commit, not after
          --- git remembers.
    \item Never echoed into a terminal, a log file, or an error
          message. Log the fact of a request, never its credentials.
    \item Never pasted into a chat transcript, including this one.
    \item Rotated if there is any doubt at all. Rotation is five
          minutes of Kim's time; a leak is permanent.
    \item \texttt{LIVE\_TRADING=0} is the master switch. The
          execution layer must refuse to place a real order unless it
          is explicitly \texttt{1}. Default closed.
\end{enumerate}
\end{resultbox}

% ══════════════════════════════════════════════════════════════════
\part{Claude's Work}

Everything from here is executable by Claude without Kim present.

\section{SSH: Getting In}

The existing instance is at \texttt{ubuntu@207.211.145.30}, with the
alias already configured in \texttt{\textasciitilde/.ssh/config}:

\begin{codebox}[~/.ssh/config]
\begin{lstlisting}[style=shell]
Host oci
  HostName 207.211.145.30
  User ubuntu
  IdentityFile ~/.ssh/oci_new
\end{lstlisting}
\end{codebox}

\begin{codebox}[Connect, and confirm what you are on]
\begin{lstlisting}[style=shell]
ssh oci                      # interactive
ssh oci 'uname -m; lsb_release -ds; python3 --version'
# expect: aarch64 / Ubuntu 24.04.4 LTS / Python 3.12.3
\end{lstlisting}
\end{codebox}

For a new instance (Option B), the key is generated locally and the
public half is handed to Oracle at instance creation:

\begin{codebox}[New host setup]
\begin{lstlisting}[style=shell]
ssh-keygen -t ed25519 -f ~/.ssh/kim_oci -C "kim-portfolio"
cat ~/.ssh/kim_oci.pub          # paste into the OCI console

cat >> ~/.ssh/config <<'EOF'

Host kim-oci
  HostName <NEW_PUBLIC_IP>
  User ubuntu
  IdentityFile ~/.ssh/kim_oci
EOF

ssh kim-oci 'echo ok'
\end{lstlisting}
\end{codebox}

\begin{resultbox}[Two things that break new OCI instances]{warn}
\begin{enumerate}[leftmargin=1.5em, topsep=2pt]
    \item \textbf{Security lists.} OCI blocks inbound traffic by
          default at the \emph{virtual network} level, independently of
          the host firewall. If SSH hangs rather than refusing, the
          ingress rule for port 22 is missing from the subnet's
          security list. Nothing on the host will fix this.
    \item \textbf{Ubuntu's iptables.} Oracle's Ubuntu images ship with
          a restrictive \texttt{iptables} ruleset that surprises people
          who have already opened the security list. Check
          \texttt{sudo iptables -L INPUT -n} before concluding the
          network is at fault.
\end{enumerate}
Outbound is open by default, which is all the scripts need.
\end{resultbox}

\section{Installing the Scripts}

One command from the Mac. The installer is idempotent --- re-running it
is the normal way to deploy a change.

\begin{codebox}[From the Mac]
\begin{lstlisting}[style=shell]
cd ~/git/andy/crypto/portfolio
bash install_oci.sh          # or: OCI_HOST=kim-oci bash install_oci.sh
\end{lstlisting}
\end{codebox}

\noindent It performs, in order:

\begin{enumerate}[leftmargin=2em]
    \item Verifies SSH reachability and fails loudly if not.
    \item Creates \texttt{/home/ubuntu/portfolio/\{reports,logs\}}.
    \item Copies \texttt{signals.py}, \texttt{weekly\_report.py},
          \texttt{run\_weekly.sh}.
    \item Builds a virtualenv and installs \texttt{numpy}, 
          \texttt{pandas}, \texttt{matplotlib}, \texttt{yfinance}.
    \item \textbf{Runs one report immediately}, so a broken install is
          discovered now rather than next Monday.
    \item Installs the weekly cron entry.
\end{enumerate}

\begin{resultbox}[ARM notes]{warn}
The instance is \texttt{aarch64}. \texttt{numpy}, \texttt{pandas} and
\texttt{matplotlib} all publish ARM wheels, so the install is a
download rather than a compile --- but it is slower than on x86 and
wants \texttt{python3-dev} present in case anything does fall back to
source. The installer handles this.

Use a virtualenv, not \texttt{pip install --user}: Ubuntu 24.04
enforces PEP~668 and will refuse to modify the system Python.

This ARM detail is also the concrete reason IBKR's Java gateway is
awkward here (Section~\ref{sec:broker}) --- its installer bundles an
x86 runtime.
\end{resultbox}

\section{The Cron Job}

The house style on this box --- visible in the ten existing entries ---
is absolute paths and appended logs. We match it.

\begin{codebox}[crontab -e]
\begin{lstlisting}[style=shell]
# Recycling portfolio weekly report
# Sunday 21:15 UTC = Monday 07:15 AEST -- before the working week,
# after the US Friday close has settled into the data.
15 21 * * 0 /home/ubuntu/portfolio/run_weekly.sh >> /home/ubuntu/portfolio/logs/cron.log 2>&1
\end{lstlisting}
\end{codebox}

The installer adds this for you, stripping any previous copy of the same
line first so repeated runs cannot produce duplicates.

\begin{resultbox}[Cron facts that cause silent failures]{warn}
\begin{enumerate}[leftmargin=1.5em, topsep=2pt]
    \item \textbf{Cron's \texttt{PATH} is nearly empty} and it does not
          read your shell profile. Every path in a cron job must be
          absolute --- including the Python interpreter. This is why
          \texttt{run\_weekly.sh} invokes
          \texttt{./venv/bin/python} rather than \texttt{python3}.
    \item \textbf{The server runs UTC.} Confirm with \texttt{date -u}
          and convert deliberately; do not assume local time.
    \item \textbf{Without redirection, output is emailed} --- and on a
          box with no mail transport, silently discarded. Always
          \texttt{>> logfile 2>\&1}.
    \item \textbf{A cron job that fails is invisible.} Check
          \texttt{grep CRON /var/log/syslog} for whether it fired at
          all, and the job's own log for what happened when it did.
    \item \textbf{Test without waiting a week:} run
          \texttt{/home/ubuntu/portfolio/run\_weekly.sh} by hand, then
          verify with \texttt{env -i /bin/bash -c
          '/home/ubuntu/portfolio/run\_weekly.sh'}, which reproduces
          cron's bare environment. Most cron bugs appear only under
          that second form.
\end{enumerate}
\end{resultbox}

\section{The Heartbeat: Not Being Reclaimed}

Oracle gives away Always Free instances, and reserves the right to take
back the ones nobody is using. A box that runs a script for four seconds
a week looks, to a monitoring system, exactly like a box nobody is
using.

\begin{resultbox}[Oracle's idle-reclamation rule, and the loophole in it]{warn}
An Always Free compute instance becomes a candidate for reclamation
when, across a \textbf{7-day window}, \textbf{all} of the following
hold:

\vspace{0.4em}
\centering
\begin{tabular}{@{}lr@{}}
    \toprule
    \textbf{Metric (95th percentile)} & \textbf{Threshold} \\
    \midrule
    CPU utilisation                     & $< 20\%$ \\
    Network utilisation                 & $< 20\%$ \\
    Memory utilisation (Ampere A1 only) & $< 20\%$ \\
    \bottomrule
\end{tabular}

\vspace{0.6em}
\raggedright
The important word is \textbf{all}. The test is a conjunction, so
\textbf{defeating any single criterion is sufficient} --- there is no
need to hold three metrics up. CPU is by far the cheapest to move, so
that is the one we move.
\end{resultbox}

\subsection{How the keepalive works}

\texttt{keepalive.sh} runs every 30 minutes and burns one core for five
minutes. That puts roughly a sixth of all samples near $100\%$ on one
core --- comfortably enough to hold the 95th percentile above $20\%$,
with a wide margin for a missed run.

It burns that core at \texttt{nice -n 19} and \texttt{ionice -c3}, the
lowest possible scheduling priority for both CPU and disk. In practice
the loop is descheduled the instant anything real wants the processor:
the four AI agents and the Gateway app never wait on it. We are
consuming \emph{idle} capacity --- which is exactly the capacity Oracle
is measuring.

\begin{codebox}[The cron pair]
\begin{lstlisting}[style=shell]
# Recycling portfolio: weekly report + keep instance from being reclaimed
15 21 * * 0   /home/ubuntu/portfolio/run_weekly.sh >> /home/ubuntu/portfolio/logs/cron.log 2>&1
*/30 * * * *  /home/ubuntu/portfolio/keepalive.sh  >> /home/ubuntu/portfolio/logs/keepalive.log 2>&1
\end{lstlisting}
\end{codebox}

\begin{resultbox}[Robin's box probably does not need this. Install it anyway.]{strong}
The existing instance runs four AI agents, the Gateway app and a
\texttt{*/15} cron job, so its CPU is unlikely to look idle. The
keepalive is insurance, and it is nearly free.

\vspace{0.4em}
It matters more than it looks, because of what happens on the day it is
needed: the agents go quiet during a lull, Oracle's 7-day window closes,
the instance is stopped, and \textbf{the weekly reports simply stop
arriving.} Nothing errors. Nothing alerts. The failure mode is
\emph{silence} --- which is precisely the failure mode a report you
glance at once a week is worst at revealing.
\end{resultbox}

\subsection{The other heartbeat: knowing it is alive}

Two different things get called a heartbeat, and we want both. The
keepalive above beats \emph{outward}, convincing Oracle we exist. We
also want one beating \emph{inward}, convincing us the job still runs.

So \texttt{keepalive.sh} writes \texttt{heartbeat.txt} every 30 minutes
\emph{before} it burns anything, so the timestamp lands even if the burn
is killed. \texttt{healthcheck.sh} reads that file, plus the report
directory and the crontab, and answers the only question that matters:

\begin{codebox}[Run this whenever you wonder]
\begin{lstlisting}[style=shell]
$ ssh oci '/home/ubuntu/portfolio/healthcheck.sh'

Portfolio health  2026-08-27T00:15:02Z

  keepalive              ok - 12 min ago
  weekly report          ok - 1 day(s) old
  report count           1 archived
  signals                no data errors
  signals                mode: paper
  cron: report           installed
  cron: keepalive        installed
  disk                   69G free of 145G (53% used)

HEALTHY
\end{lstlisting}
\end{codebox}

It exits non-zero when anything needs attention, so it composes with
anything that watches exit codes. The checks that fire in practice: a
stale keepalive (instance stopped, or cron broken), an overdue report
(more than 8 days old), and data-feed errors recorded by the signal
engine.

\begin{resultbox}[Free-tier hygiene, restated because it is the only thing that costs money]{fail}
Compute on the Always Free tier is free and stays free.
\textbf{Storage is capped at 200~GB across the entire tenancy ---
including volumes attached to nothing at all.}

On 3 June 2026 four abandoned 47~GB rescue boot volumes pushed Robin's
tenancy to 338~GB and produced roughly AUD~\$8/month of charges on what
should have been a \$0 bill. \texttt{healthcheck.sh} reports disk usage
for this reason. If a bill ever appears, look for orphaned volumes
first --- the commands are in Section~\ref{sec:manifest}.
\end{resultbox}

\section{The Signal Engine}

\texttt{signals.py} computes each asset's momentum ratio with arithmetic
\emph{identical} to \texttt{momentum\_single.py}, so live signals and
backtested signals cannot silently diverge. It then applies the
recycling rule with the Solana cap.

\begin{codebox}[Verified output, run 2026-08-27]
\begin{lstlisting}[style=shell]
$ ./venv/bin/python signals.py

Signals as of 2026-08-26T21:32:07+00:00  (3/4 active)

Asset    Symbol           Close    Ratio  Signal   Weight  Trade as
----------------------------------------------------------------------
COPPER   HG=F              6.70    0.817    cash     0.0%  CPER @ IBKR
NATGAS   NG=F              2.90    1.102    HOLD    37.5%  UNG @ IBKR
BTC      BTC-USD      78,659.15    1.184    HOLD    37.5%  XBT/USD @ KRAKEN
SOL      SOL-USD          97.59    1.302    HOLD    25.0%  SOL/USD @ KRAKEN (capped)
----------------------------------------------------------------------
CASH                                                 0.0%
\end{lstlisting}
\end{codebox}

Note that the cap is binding in this very sample: three assets are
active, so equal weighting would give Solana $33.3\%$, and the cap
holds it to $25\%$ with the excess going to copper and gas. The
mechanism is not hypothetical.

\subsection{Two deliberate design decisions}

\begin{enumerate}[leftmargin=2em]
    \item \textbf{Fail-safe on missing data.} If a price feed fails, the
          asset's signal is recorded as \emph{off} and the money stays
          in cash. A data outage can therefore cost us an opportunity;
          it can never cause an accidental position. Errors are listed
          in the report so a silent outage cannot masquerade as a
          genuine ``all clear.''
    \item \textbf{Signal instrument $\neq$ traded instrument.} Momentum
          is computed on \texttt{HG=F} and \texttt{NG=F} --- exactly what
          the study measured --- while orders would be placed in
          \texttt{CPER}/\texttt{UNG}. Conflating the two would quietly
          invalidate the research, so the mapping is explicit in the
          \texttt{ASSETS} table.
\end{enumerate}

\section{The Execution Layer (Not Yet Written)}
\label{sec:exec}

This is the only substantial piece of code still outstanding, and it
should not be written until Phase 2 has run clean. The specification:

\begin{enumerate}[leftmargin=2em]
    \item Read target weights from \texttt{signals.py}.
    \item Read actual holdings from each venue.
    \item Compute the difference, and \textbf{only trade if a position
          is off target by more than a threshold} (5\% of portfolio
          value is a sensible starting point). Without this, tiny daily
          drift generates constant fee-burning churn.
    \item Place market orders, sized to the smaller of the required
          trade and \texttt{MAX\_ORDER\_USD}.
    \item Refuse to do anything at all unless
          \texttt{LIVE\_TRADING=1}. Log the orders it \emph{would} have
          placed otherwise --- that log is how we validate it.
    \item Write \texttt{accounts.json} after every run, which is what
          the weekly report reads to reconcile paper against real.
\end{enumerate}

\begin{codebox}[Sketch: reading balances from Kraken via ccxt]
\begin{lstlisting}[style=py]
import os, ccxt

kraken = ccxt.kraken({
    "apiKey": os.environ["KRAKEN_API_KEY"],
    "secret": os.environ["KRAKEN_API_SECRET"],
    "enableRateLimit": True,
})

balance = kraken.fetch_balance()          # needs "Query funds" only
print({k: v for k, v in balance["total"].items() if v})

# Placing an order needs "Create & modify orders". Guarded:
if os.environ.get("LIVE_TRADING") != "1":
    raise SystemExit("LIVE_TRADING is not 1 -- refusing to place orders")
\end{lstlisting}
\end{codebox}

\begin{resultbox}[Build it in this order, and not another]{strong}
\begin{enumerate}[leftmargin=1.5em, topsep=2pt]
    \item Read-only first. Fetch balances and prices for a week.
          Confirm the numbers match what Kim sees when he logs in.
    \item Then dry-run orders --- computed, logged, never sent ---
          for a week. Read every line.
    \item Then live orders capped at \$50, for a week.
    \item Then raise the cap.
\end{enumerate}
Each stage catches a category of bug the next stage would make
expensive. Skipping stage 1 is how people discover their position
sizing was wrong by a factor of a hundred.
\end{resultbox}

% ══════════════════════════════════════════════════════════════════
\part{The Weekly Report}

\section{What It Produces}

Every Monday morning the server writes a two-page PDF and a JSON
summary. The PDF is rendered with \texttt{matplotlib}, deliberately
\emph{not} LaTeX --- the server needs no TeX installation, and there is
no compile step that can fail at 7am.

\textbf{Page 1:} headline tiles (NAV, weekly change, return since
inception, maximum drawdown); the signal table with each asset's ratio,
target weight and traded instrument; a bar showing where the money
should be right now; and a notes block flagging paper-versus-live mode,
data errors, and the Solana cap.

\textbf{Page 2:} the NAV curve since inception, the drawdown curve, and
a timeline strip showing which assets were held when.

\begin{codebox}[Verified output, run 2026-08-27]
\begin{lstlisting}[style=shell]
$ ./venv/bin/python weekly_report.py --outdir ~/portfolio/reports
Wrote /home/ubuntu/portfolio/reports/portfolio-2026-W35.pdf
  Paper NAV $11,862.83  week +10.81%  total +18.63%  maxDD -2.9%
  Mode: paper   Active: 3/4
\end{lstlisting}
\end{codebox}

\section{The Design Decision That Matters}

\begin{resultbox}[The report is stateless]{strong}
Every run rebuilds the \emph{entire} equity curve from scratch, from
first principles, back to inception. Nothing is carried forward
between runs.

This costs a few seconds of ARM CPU per week, and buys a property
worth far more than that: \textbf{a missed week, a crashed run, a
truncated file, or a server rebuild cannot cause the tracked NAV to
drift.} An incrementally-updated ledger accumulates errors silently
and is only discovered when the number is already wrong. A recomputed
one is self-healing --- correct this week regardless of what happened
last week.

The same reasoning is why \texttt{run\_weekly.sh} keeps the last 52
reports rather than overwriting: the history is evidence, and
evidence should be append-only.
\end{resultbox}

\section{How Claude Fetches It}

Three ways, in increasing order of convenience:

\begin{codebox}[Getting at the report]
\begin{lstlisting}[style=shell]
# 1. The one-liner Claude will normally use: numbers, no file transfer.
ssh oci 'cat /home/ubuntu/portfolio/reports/latest.json'

# 2. Copy the PDF back and open it.
bash ~/git/andy/crypto/portfolio/fetch_report.sh          # latest
bash ~/git/andy/crypto/portfolio/fetch_report.sh 2026-W35 # a given week

# 3. Look around.
ssh oci 'ls -lt /home/ubuntu/portfolio/reports/ | head'
ssh oci 'tail -40 /home/ubuntu/portfolio/logs/cron.log'
\end{lstlisting}
\end{codebox}

\texttt{latest.json} exists precisely so Claude does not have to parse a
PDF to answer ``how is it going?''. It carries the week, the mode
(paper or live), every ratio, every weight, and the summary statistics.
\texttt{fetch\_report.sh} prints a one-line summary and opens the PDF in
Skim.

\begin{codebox}[latest.json, abridged]
\begin{lstlisting}[style=shell]
{
  "week": "2026-W35",
  "mode": "paper",
  "stats": { "nav": 11862.83, "week_return": 0.1081,
             "total_return": 0.1863, "max_dd": -0.029, "sharpe": ... },
  "signals": { "COPPER": false, "NATGAS": true, "BTC": true, "SOL": true },
  "ratios":  { "COPPER": 0.817, "NATGAS": 1.102, "BTC": 1.184, "SOL": 1.302 },
  "weights": { "COPPER": 0.0, "NATGAS": 0.375, "BTC": 0.375, "SOL": 0.25 },
  "cash": 0.0, "errors": []
}
\end{lstlisting}
\end{codebox}

% ══════════════════════════════════════════════════════════════════
\part{Going Live}

\section{Choosing the Solana Cap}
\label{sec:cap}

This was an open question when the runbook was drafted, and it is now
closed. \texttt{cap\_sweep.py} regenerates momentum for all four assets
from source prices --- the saved \texttt{timing\_*.csv} files could not
be used, see the warning below --- builds the recycling portfolio for a
range of caps, and reports the trade-off.

Period 19 April 2021 to 26 August 2026 (1{,}956 days, 5.35 years),
\$10{,}000 start, signals lagged one day, no fees.

\vspace{0.5em}
\centering
{\small
\begin{tabular}{@{}lrrrrr@{}}
    \toprule
    \textbf{SOL cap} & \textbf{Final} & \textbf{CAGR} & \textbf{Sharpe} & \textbf{Sortino} & \textbf{Max DD} \\
    \midrule
    $0\%$ (no SOL)   & \$95{,}575  & $52.5\%$ & $1.26$ & $1.43$ & $-35.3\%$ \\
    $10\%$           & \$119{,}632 & $59.0\%$ & $1.37$ & $1.71$ & $-35.1\%$ \\
    $15\%$           & \$132{,}434 & $62.0\%$ & $1.41$ & $1.79$ & $-35.2\%$ \\
    $17.5\%$         & \$138{,}973 & $63.5\%$ & $1.42$ & $1.82$ & $-35.3\%$ \\
    $20\%$           & \$145{,}579 & $64.9\%$ & $1.43$ & $1.86$ & $-36.3\%$ \\
    \textbf{$25\%$}  & \textbf{\$158{,}908} & \textbf{$67.7\%$} & \textbf{$1.44$} & \textbf{$1.91$} & \textbf{$-38.7\%$} \\
    $33.3\%$         & \$178{,}679 & $71.4\%$ & $1.44$ & $1.94$ & $-42.8\%$ \\
    $50\%$           & \$169{,}259 & $69.6\%$ & $1.35$ & $1.81$ & $-50.6\%$ \\
    \textcolor{red}{no cap} & \textcolor{red}{\$114{,}144} & \textcolor{red}{$57.6\%$}
                     & \textcolor{red}{$1.08$} & \textcolor{red}{$1.37$} & \textcolor{red}{$-75.4\%$} \\
    \bottomrule
\end{tabular}
}

\raggedright
\vspace{0.8em}

\begin{resultbox}[The result nobody expected: the cap is not a trade-off]{strong}
\textbf{The uncapped rule is the worst row in the table.} Not the
riskiest-but-highest-returning --- the worst, on both axes at once. It
makes less money than a $10\%$ cap and takes double the drawdown.
\emph{Every} cap between $10\%$ and $50\%$ dominates it.

\vspace{0.5em}
The mechanism is volatility drag. Solana signals alone on
\textbf{$24.1\%$ of its active days}, and on those days the uncapped
rule puts the entire portfolio into an asset that has historically
fallen $96\%$. Occasionally concentrating everything into that
destroys the compounding, however good the asset is on average. This is
the $-50\%$-then-$+50\%$ arithmetic from \texttt{concepts.pdf},
operating at portfolio scale.

\vspace{0.5em}
So the cap does not buy safety at the cost of return. \textbf{It buys
both.}
\end{resultbox}

\subsection{Why 25\%, given the plateau}

Sharpe is essentially flat from $20\%$ to $33\%$ ($1.43$ to $1.45$) ---
those differences are noise, not signal. Return keeps climbing to
$33\%$; drawdown climbs with it, and faster beyond $25\%$. So the choice
inside the plateau is a judgement, and it should not be made by picking
the maximum of anything.

\vspace{0.5em}
A split-half check is the reason for caution:

\vspace{0.5em}
\centering
\begin{tabular}{@{}lll@{}}
    \toprule
    & \textbf{First half (to Dec 2023)} & \textbf{Second half (to Aug 2026)} \\
    \midrule
    Best Sharpe & $33\%$ & $15\%$ \\
    Best Calmar & $33\%$ & $27.5\%$ \\
    Best return & $33\%$ & $50\%$ \\
    Uncapped    & disastrous ($-75\%$ DD, Sharpe $0.83$) & barely differs from capped \\
    \bottomrule
\end{tabular}

\raggedright
\vspace{0.8em}

The two halves disagree about the optimum and agree about everything
that matters: caps in the $15$--$33\%$ band beat both extremes, and
uncapped is only catastrophic in the half that contained a crash. A
number tuned to the full-period maximum would be fitting the 2022
drawdown.

\begin{resultbox}[Decision: 25\%]{strong}
\begin{enumerate}[leftmargin=1.5em, topsep=2pt]
    \item \textbf{It is not a fitted parameter.} $25\%$ is the natural
          equal share of four assets. The cap therefore only ever binds
          \emph{downward} --- it can never force us to hold more Solana
          than plain diversification would. That is the strongest
          available defence against overfitting: the number comes from
          the portfolio's structure, not from the data.
    \item It sits inside the flat region of the Sharpe curve, so it
          gives up nothing measurable in risk-adjusted terms.
    \item It adds $66\%$ to the final value over the three-asset
          portfolio (\$95{,}575 $\rightarrow$ \$158{,}908) while moving
          maximum drawdown only from $-35.3\%$ to $-38.7\%$.
    \item Both halves of the sample support it.
\end{enumerate}

\vspace{0.5em}
\textbf{If Kim wants a gentler ride,} $20\%$ is the defensible
alternative: the best Calmar ratio in the table, drawdown a little over
a point shallower, at the cost of about $8\%$ of the final value. That
is a preference, not a correction --- both are right.

\vspace{0.4em}
\textbf{Do not raise it above $33\%$.} Past that, drawdown grows and
return stops improving. $\texttt{SOL\_CAP}$ in \texttt{.env} exists so
this is one line to change, and \texttt{signals.py} enforces it on
every run.
\end{resultbox}

\begin{resultbox}[The data problem this uncovered]{fail}
The saved \texttt{timing\_*.csv} files in
\texttt{\textasciitilde/git/andy/crypto/} \textbf{are stale and
inconsistent with the published study}, and any analysis built on them
would have been confidently wrong:

\vspace{0.4em}
\centering
\begin{tabular}{@{}llll@{}}
    \toprule
    \textbf{File} & \textbf{Contains} & \textbf{Study says} & \\
    \midrule
    \texttt{timing\_BTCUSD.csv} & $+462.8\%$ & $+462.7\%$ & \textcolor{green}{matches} \\
    \texttt{timing\_SOLUSD.csv} & $+201.8\%$ & $+7{,}691\%$ & \textcolor{red}{365-day run} \\
    \texttt{timing\_HG=F.csv}   & $+15.0\%$  & $+97.0\%$    & \textcolor{red}{wrong run} \\
    \texttt{timing\_NG=F.csv}   & $+88.4\%$  & $+309.9\%$   & \textcolor{red}{wrong run} \\
    \bottomrule
\end{tabular}

\vspace{0.5em}
\raggedright
The Solana file holds the \textbf{365-day} lookback --- the one variant
that produced \emph{negative} alpha --- not the 180-day run the study
recommends. Copper and gas match neither.

\vspace{0.4em}
\texttt{cap\_sweep.py} therefore ignores those files entirely and
refetches prices into its own \texttt{cache/} directory. Leftover
result files from a parameter sweep are not a data source. They should
be regenerated or deleted before they mislead someone.
\end{resultbox}

\section{Four Phases, With Gates}

Each phase has an exit criterion. The point of a gate is that it can
fail --- if a criterion is not met, the answer is to stop, not to
proceed while noting a concern.

\vspace{0.5em}
\begin{tabular}{@{}p{1.5cm}p{3.1cm}p{4.0cm}p{4.6cm}@{}}
    \toprule
    \textbf{Phase} & \textbf{What runs} & \textbf{Kim does} & \textbf{Gate to pass} \\
    \midrule
    \textbf{1}\newline 1 week &
    Reports on Robin's box, paper only, no keys &
    Step 1 (Oracle); nothing else &
    A report lands unattended on Monday; signals match a hand-check \\[0.4em]
    \textbf{2}\newline 4--12 weeks &
    Execution layer dry-run; ETF tracking measured vs futures &
    Steps 2--3 (accounts, paper keys) &
    Dry-run orders sane every week; backtest confirms the cap; ETF decay measured \\[0.4em]
    \textbf{3}\newline 4 weeks &
    Live crypto leg only, \$50 order cap &
    Step 4 (fund small), live keys &
    Broker NAV tracks paper NAV within a few percent; no failed orders \\[0.4em]
    \textbf{4}\newline ongoing &
    Full size; copper/gas added if Phase 2 justified it &
    Increase funding; hardware wallet &
    Reviewed quarterly against the mind-changers below \\
    \bottomrule
\end{tabular}

\vspace{0.8em}
\begin{resultbox}[The gate on Phase 2 is the one that matters]{warn}
One of the two gaps here is now closed. The four-asset backtest with the
Solana cap has been run --- Section~\ref{sec:cap} --- and it supports the
design: capped at $25\%$, \$10{,}000 becomes \$158{,}908 over 5.35 years
with a $-38.7\%$ worst drawdown, against \$95{,}575 for the three-asset
portfolio. Uncapped is worse on \emph{both} counts.

\vspace{0.5em}
What Phase 2 still has to establish is the part no backtest can tell us:
\textbf{how much of that survives in instruments we can actually buy.}
The signals are measured on \texttt{HG=F} and \texttt{NG=F}; the orders
would be in \texttt{CPER} and \texttt{UNG}, which carry roll decay the
backtest never saw. Until that difference is measured on live prices,
the copper and gas legs stay on paper.
\end{resultbox}

\section{What Would Stop Us}

Stated in advance, so they cannot be rationalised away later:

\begin{itemize}[leftmargin=2em]
    \item ETF tracking error eats the copper or gas alpha
          $\Rightarrow$ drop that leg rather than trade it hoping.
    \item Signal overlap between the four assets turns out to be high
          $\Rightarrow$ the recycling rationale is void; rebuild the
          allocation logic.
    \item UNG's roll decay eats the natural gas alpha $\Rightarrow$ drop
          the gas leg or find a different instrument. Do not trade it
          hoping.
    \item Live results diverge materially from paper over a quarter
          $\Rightarrow$ stop and find out why before adding money.
    \item Results prove sensitive to the lookback window $\Rightarrow$
          treat the whole thing as overfitted. A real effect should not
          require the knob set to exactly 180.
    \item Kim cannot comfortably tolerate a 30\% drawdown in practice
          rather than in principle $\Rightarrow$ reduce size until he
          can. This is a legitimate reason and needs no other
          justification.
\end{itemize}

\section{Command Reference}
\label{sec:manifest}

\begin{codebox}[Everything, in one place]
\begin{lstlisting}[style=shell]
# --- Deploy (from the Mac) -------------------------------------------
cd ~/git/andy/crypto/portfolio
bash install_oci.sh                    # install + venv + cron + smoke test
OCI_HOST=kim-oci bash install_oci.sh   # ... to a different host

# --- Check on it ------------------------------------------------------
ssh oci 'cat /home/ubuntu/portfolio/reports/latest.json'
ssh oci 'cd /home/ubuntu/portfolio && ./venv/bin/python signals.py'
ssh oci 'tail -40 /home/ubuntu/portfolio/logs/cron.log'
ssh oci 'crontab -l | grep -A2 Recycling'
ssh oci '/home/ubuntu/portfolio/healthcheck.sh'     # alive and on schedule?
ssh oci 'cat /home/ubuntu/portfolio/heartbeat.txt'  # last keepalive
ssh oci 'grep CRON /var/log/syslog | tail -5'

# --- Reports ----------------------------------------------------------
bash fetch_report.sh                   # newest, opens in Skim
bash fetch_report.sh 2026-W35          # a specific week
ssh oci 'ls -lt /home/ubuntu/portfolio/reports | head'

# --- Run out of band --------------------------------------------------
ssh oci '/home/ubuntu/portfolio/run_weekly.sh'
ssh oci "env -i /bin/bash -c '/home/ubuntu/portfolio/run_weekly.sh'"  # as cron sees it

# --- Free-tier hygiene (do this after any OCI console work) -----------
# Orphaned volumes are the only way this costs money. 200 GB total cap.
oci bv boot-volume list --compartment-id <tenancy-ocid> \
    --availability-domain "drxR:AP-MELBOURNE-1-AD-1" --all
oci bv volume list --compartment-id <tenancy-ocid> --all
\end{lstlisting}
\end{codebox}

\subsection*{File manifest}

\vspace{0.3em}
\begin{tabular}{@{}p{4.6cm}p{9.3cm}@{}}
    \toprule
    \textbf{File} & \textbf{Purpose} \\
    \midrule
    \texttt{signals.py} &
    Momentum ratios for the four assets; recycling weights with the
    Solana cap. Arithmetic identical to \texttt{momentum\_single.py}. \\
    \texttt{weekly\_report.py} &
    Rebuilds the paper equity curve from scratch; renders the two-page
    PDF, \texttt{latest.pdf} and \texttt{latest.json}. \\
    \texttt{run\_weekly.sh} &
    Cron entry point. Absolute paths, logs, prunes to 52 reports. \\
    \texttt{install\_oci.sh} &
    Idempotent deploy from the Mac: copy, venv, smoke test, cron. \\
    \texttt{fetch\_report.sh} &
    Pull the latest PDF back to the Mac and open it. \\
    \texttt{keepalive.sh} &
    Burns one core for 5 min every 30 min so Oracle does not reclaim the
    instance; writes \texttt{heartbeat.txt}. \\
    \texttt{healthcheck.sh} &
    One command answering ``is it alive and on schedule?''. Non-zero exit
    if not. \\
    \texttt{cap\_sweep.py} &
    Regenerates momentum from source data and sweeps the Solana cap.
    Produced the numbers in Section~\ref{sec:cap}. \\
    \texttt{make\_figure.py} &
    Self-contained: fetches prices and draws the three-way portfolio
    comparison. Listed in full in \texttt{algorithm.pdf}. \\
    \texttt{.env} (server only) &
    API keys and risk limits. \texttt{chmod 600}, never in git. \\
    \texttt{execute.py} &
    \textcolor{red}{Not yet written.} Order placement. Phase 2. \\
    \bottomrule
\end{tabular}

\vspace{1.5em}
\begin{resultbox}[]{warn}
\begin{center}
\Large\bfseries Nothing in this document is financial advice.
\end{center}
\vspace{0.3em}
\centering
It is a deployment runbook for a strategy that has been backtested and
not yet validated live. The strategy's own caveats are in
\texttt{concepts.pdf} and \texttt{kim.pdf}, and they are substantial.
Read them before Step 4.
\end{resultbox}


% ══════════════════════════════════════════════════════════════════
\newpage
\part{Appendix: The Complete Code}

Everything Claude needs to implement the copper--gas--bitcoin--sol
strategy is listed here in full. Nothing is elided.

\begin{resultbox}[These listings cannot go stale]{strong}
The listings below are pulled from the working files at compile time
via \texttt{\textbackslash lstinputlisting}. They are not copies. If a
script changes and this document is rebuilt, the appendix changes with
it --- so the PDF and the code can never silently disagree.

\vspace{0.5em}
Source of truth:
\texttt{\textasciitilde/git/andy/crypto/portfolio/}.
Recompile with \texttt{tectonic oci.tex} after any edit.
\end{resultbox}

\vspace{0.5em}
\begin{tabular}{@{}p{3.6cm}rp{9.4cm}@{}}
    \toprule
    \textbf{File} & \textbf{Lines} & \textbf{Role} \\
    \midrule
    \texttt{signals.py}        & 229 & Momentum ratios; recycling weights with the SOL cap \\
    \texttt{weekly\_report.py} & 413 & Paper equity curve; the two-page PDF report \\
    \texttt{cap\_sweep.py}     & 213 & Cap selection backtest (Section~\ref{sec:cap}) \\
    \texttt{run\_weekly.sh}    & 14  & Cron entry point \\
    \texttt{keepalive.sh}      & 47  & Oracle idle-reclamation defence + heartbeat \\
    \texttt{healthcheck.sh}    & 83  & Liveness and schedule check \\
    \texttt{install\_oci.sh}   & 72  & One-command deploy from the Mac \\
    \texttt{fetch\_report.sh}  & 25  & Pull the latest report back \\
    \midrule
    \textbf{Total}             & \textbf{1{,}096} & \\
    \bottomrule
\end{tabular}

\vspace{0.8em}
\begin{resultbox}[What is deliberately NOT here]{fail}
\texttt{execute.py} --- the code that places real orders --- is not in
this appendix because it has not been written, and it should not be
written until Phase~2 has run clean. Its specification is in
Section~\ref{sec:exec}, and the build order given there
(read-only $\rightarrow$ dry-run $\rightarrow$ \$50 cap $\rightarrow$
full size) matters more than the code itself.

\vspace{0.4em}
Also absent: \texttt{.env}, which holds the API keys. It is created on
the server by hand, \texttt{chmod 600}, and never enters git or any
document. See Section~\ref{sec:keys}.
\end{resultbox}

% ──────────────────────────────────────────────────────────────────
\section{signals.py}

The heart of the system. The momentum arithmetic here is byte-for-byte
equivalent to \texttt{momentum\_single.py}, which is what makes live
signals comparable with the backtest. Note two deliberate choices: a
missing price feed produces \emph{cash}, never a position; and the
signal instrument is kept distinct from the traded instrument.

\lstinputlisting[style=py]{portfolio/signals.py}

% ──────────────────────────────────────────────────────────────────
\section{weekly\_report.py}

Rebuilds the entire paper equity curve from scratch on every run, then
renders the report with \texttt{matplotlib} --- no LaTeX on the server,
so there is no compile step that can fail at 7am.

\lstinputlisting[style=py]{portfolio/weekly_report.py}

% ──────────────────────────────────────────────────────────────────
\section{cap\_sweep.py}

Produced the cap decision in Section~\ref{sec:cap}. It refetches prices
into its own \texttt{cache/} rather than trusting the stale
\texttt{timing\_*.csv} files, and handles the crypto/futures calendar
mismatch by treating the 7-day crypto calendar as master and
forward-filling futures closes.

\lstinputlisting[style=py]{portfolio/cap_sweep.py}

% ──────────────────────────────────────────────────────────────────
\section{run\_weekly.sh}

The cron entry point. Absolute paths throughout, because cron's
\texttt{PATH} is nearly empty and it does not read a shell profile.

\lstinputlisting[style=shell]{portfolio/run_weekly.sh}

% ──────────────────────────────────────────────────────────────────
\section{keepalive.sh}

Stops Oracle reclaiming the instance, and writes the heartbeat used by
\texttt{healthcheck.sh}. The heartbeat is written \emph{before} the CPU
burn so the timestamp survives even if the burn is killed.

\lstinputlisting[style=shell]{portfolio/keepalive.sh}

% ──────────────────────────────────────────────────────────────────
\section{healthcheck.sh}

Answers ``is it alive and on schedule?'' and exits non-zero if not.

\lstinputlisting[style=shell]{portfolio/healthcheck.sh}

% ──────────────────────────────────────────────────────────────────
\section{install\_oci.sh}

Idempotent deploy from the Mac. Re-running it is the normal way to ship
a change. It runs one report immediately, so a broken install is
discovered now rather than next Monday.

\lstinputlisting[style=shell]{portfolio/install_oci.sh}

% ──────────────────────────────────────────────────────────────────
\section{fetch\_report.sh}

Pulls the latest report back to the Mac and opens it in Skim.

\lstinputlisting[style=shell]{portfolio/fetch_report.sh}

% ------------------------------------------------------------------
\section{make\_figure.py}

Self-contained: fetches prices from Yahoo Finance and draws the
three-way portfolio comparison with no dependency on any other file
here. Presented with a walkthrough in \texttt{algorithm.pdf}.

\lstinputlisting[style=py]{portfolio/make_figure.py}

\end{document}
